\section{Хеширование строк}

\subsection{Домашнее задание}

  \begin{definition} Семейство хеш-функций $\H = \{ h : X \rightarrow Y \}$ называется
    \begin{itemize}
      \item универсальным, если:
        $$\forall x_1, x_2 \in X, x_1 \neq x_2: \Prb{h \in \H}\left[ \; h(x_1) = h(x_2) \; \right] \le \frac{1}{|Y|}$$

      \item $k$-независимым, если 
        для любых различных $x_1, x_2, \cdots, x_k \in X$, для любых, возможно, совпадающих $y_1, y_2, \cdots, y_k \in Y$ выполняется:
        $$\Prb{h \in \H}\left[ \; \bigwedge_{i=1}^k h(x_i) = y_i \; \right] = \frac{1}{|Y|^k}$$
    \end{itemize}
  \end{definition}

\begin{enumerate}

  \item
    \begin{enumerate}
	  \item Докажите, что любое 2-независимое семейство хэш-функций является универсальным.
	  \item Докажите, что любое $k+1$-независимое семейство хэш-функций является $k$-независимым.
    \end{enumerate}

  \item
    Придумайте строку длины $n$ над алфавитом $\{0, 1\}$, в которой $\Omega(n^2)$ различных подстрок.

  \item
    Даны две строки суммарной длины $n$. Найдите их самую длинную общую подстроку. $\O(n \log n)$.

  \item
    Даны строка $s$ длины $n$ и число $k$. Найдите в $s$ за время $\O(n)$ самую длинную подстроку,
    представимую как степень  некоторой строки $x$, для которой $|x| = k$ (степенью строки
    называют её конкатенацию с собой несколько раз).

  \item
    Найти подстроку в тексте. При сравнении строк, если несовпадений символов было не более~$k$, строки считаются равными. $\O(nk \log n)$.

  \item
    Дана строка длины $n$. Запросы:
    \begin{itemize}
      \item \texttt{set(i, c)} -- установить символ на $i$-й позиции в $c$.
      \item \texttt{test(i, j, l)} -- проверить на равенство подстроки длины $l$ начинающиеся в $i$-й и $j$-й позициях.
    \end{itemize}
    $\O(\log{n})$ на запрос, online.

\end{enumerate}